\chapter{Variables aléatoires et espérance}
\section*{Exercice 21.1 -- Établir une loi}
On lance deux pièces équilibrées indépendantes. Le gain brut vaut $4$ pour deux faces, $1$ pour une seule face et $0$ sinon. Donner sa loi et son espérance.

\section*{Exercice 21.2 -- Mise équitable}
Un jeu verse $12$ avec probabilité $1/6$, $3$ avec probabilité $1/3$ et $0$ sinon. Quelle mise rend le gain net d'espérance nulle ?

\section*{Exercice 21.3 -- Linéarité sans indépendance}
Si $X$ est le résultat d'un dé équilibré et $Y=7-X$, calculer $\mathbb E(X+Y)$ de deux façons.

\section*{Exercice 21.4 -- Compter avec des indicatrices}
Pour quatre événements de probabilité $0{,}3$, calculer l'espérance du nombre d'événements réalisés. Faut-il l'indépendance ?

\section*{Exercice 21.5 -- Transformation non linéaire}
$X$ vaut $0$ ou $2$ avec probabilités égales. Comparer $\mathbb E(X^2)$ et $\mathbb E(X)^2$.

\section*{Exercice 21.6 -- Bornes de l'espérance}
Si $a\le X\le b$ presque sûrement, démontrer $a\le\mathbb E(X)\le b$. Quand peut-on avoir $\mathbb E(X)=a$ ?
